\chapter{Ingeniería de requisitos}

\section{Introducción a la ingeniería de requisitos}

Determinar qué debe hacer un sistema es más difícil que construirlo. El problema
no es técnico: consiste en enunciar con precisión un problema complejo del que
las personas implicadas tienen ideas distintas, a veces incompatibles y casi
siempre incompletas.

Un \textbf{requisito} es una condición o capacidad que el sistema debe satisfacer.
La palabra se usa con dos niveles de detalle que conviene no mezclar
\cite{sommerville2016}:

\begin{itemize}
    \item Los \textbf{requisitos de usuario} se enuncian en lenguaje natural,
          desde el punto de vista de quien va a usar el sistema, y sirven para
          acordar el alcance con el cliente.
    \item Los \textbf{requisitos del sistema} detallan las funciones, servicios y
          restricciones con la precisión necesaria para diseñar y contratar.
\end{itemize}

La distinción clásica es otra, y es transversal a la anterior:

\begin{itemize}
    \item \textbf{Requisitos funcionales}: lo que el sistema debe hacer.
    \item \textbf{Requisitos no funcionales}: las restricciones sobre cómo debe
          hacerlo —rendimiento, disponibilidad, seguridad, normativa aplicable—.
\end{itemize}

Los no funcionales suelen recibir menos atención y son los que más condicionan la
arquitectura. Un requisito de tiempo de respuesta o de tolerancia a fallos afecta
al sistema entero, mientras que un requisito funcional suele quedar confinado a
un módulo. Además, incumplir un requisito no funcional puede inutilizar el sistema
aunque todas sus funciones estén implementadas.

\subsection*{Propiedades de un buen conjunto de requisitos}

Un requisito bien escrito es \emph{necesario}, \emph{verificable}, \emph{no
ambiguo} y \emph{consistente} con los demás. La verificabilidad es la propiedad
que más se descuida: «el sistema debe ser rápido» no es un requisito, porque no
existe prueba que decida si se cumple. «El sistema debe responder a una consulta
de catálogo en menos de dos segundos con cien usuarios concurrentes» sí lo es.

Al conjunto se le exige además que sea \emph{completo} y \emph{trazable}: que
cubra todo lo necesario y que cada elemento del diseño y de las pruebas pueda
remontarse al requisito que lo justifica. La completitud absoluta es inalcanzable
en sistemas grandes, y esa es una de las razones por las que los procesos
iterativos ganaron terreno.

\subsection*{Delimitar lo que queda fuera}

Tan importante como decir qué hace el sistema es dejar por escrito qué no hace. Un
alcance que no se cierra crece durante el desarrollo sin que nadie lo decida, y
ese crecimiento es una de las causas habituales de que un proyecto se desvíe en
plazo y coste. La lista de exclusiones debe ser explícita y estar acordada.

\section{Obtención de requisitos}

La obtención —o elicitación— es la actividad de descubrir los requisitos hablando
con las personas implicadas y observando cómo trabajan. Las dificultades son
recurrentes: quienes van a usar el sistema describen su trabajo en términos de la
herramienta que ya conocen, dan por supuesto lo que para ellos es evidente, y
distintos grupos tienen objetivos que chocan.

Las técnicas más usadas son:

\begin{itemize}
    \item \textbf{Entrevistas}, abiertas o estructuradas. Son el instrumento
          principal. Las abiertas sirven para explorar; las estructuradas, para
          confirmar y precisar.
    \item \textbf{Observación} del trabajo real. Descubre lo que nadie menciona en
          una entrevista porque lo hace de forma automática, y también los atajos
          con los que se sortea el sistema actual.
    \item \textbf{Talleres} con varios grupos a la vez, útiles para sacar a la luz
          los conflictos entre ellos en vez de descubrirlos al final.
    \item \textbf{Prototipos} y maquetas, que convierten una discusión abstracta
          en una reacción concreta ante algo visible.
    \item \textbf{Estudio de la documentación} existente y del sistema al que se
          sustituye.
\end{itemize}

Identificar a los \emph{stakeholders} —toda persona u organización afectada por el
sistema— forma parte de la actividad. Olvidar a un grupo es la vía más directa a
un requisito descubierto tarde: los administradores del sistema, el personal de
soporte y los responsables de cumplimiento normativo quedan fuera de la lista con
frecuencia y aportan restricciones que luego resultan innegociables.

\section{Modelado de casos de uso e historias de usuario}

\subsection*{Casos de uso}

Un caso de uso describe una interacción entre el sistema y un actor externo que
produce un resultado con valor para ese actor. El \textbf{actor} es un rol, no una
persona: el mismo empleado puede actuar como «cliente» y como «administrador» en
casos distintos.

El diagrama de casos de uso de UML solo aporta la vista general: qué actores hay y
con qué casos se relacionan. El contenido está en la descripción textual de cada
caso, que suele incluir el actor principal, las precondiciones, el flujo
principal paso a paso, los flujos alternativos y la postcondición
\cite{arlow2006,larman2006}.

Las relaciones entre casos —\texttt{include} para el comportamiento común y
\texttt{extend} para el opcional— conviene usarlas con moderación. Un modelo lleno
de descomposiciones se vuelve ilegible y empieza a describir el diseño en vez del
comportamiento observable, que es justo lo que un caso de uso debe evitar.

El error más frecuente al modelar es escribir los casos desde dentro del sistema,
en términos de pantallas y de operaciones sobre la base de datos. Un caso de uso
debe poder leerlo el cliente y reconocer su trabajo en él.

\subsection*{Historias de usuario}

Las historias de usuario son la alternativa ágil: enunciados breves, en la forma
«como \emph{rol} quiero \emph{acción} para \emph{beneficio}», acompañados de
criterios de aceptación. No pretenden documentar el sistema, sino recordar una
conversación pendiente con el cliente y servir de unidad de planificación.

Frente a los casos de uso son más ligeras y más fáciles de reordenar por
prioridad, pero dejan menos rastro. La elección entre unos y otras no es de moda:
depende de si el cliente está disponible de forma continua —caso en el que la
conversación sustituye al documento— o de si el acuerdo tiene que sostenerse por
sí solo durante meses.

\section{Especificación y análisis}

\subsection*{Especificación}

Especificar es dejar los requisitos por escrito de forma que sirvan de contrato
técnico. El documento de especificación —el \emph{SRS}, por sus siglas en
inglés— recoge la introducción y el alcance, la descripción general del sistema y
su entorno, los requisitos funcionales, los no funcionales y las restricciones,
más los apéndices con glosario y modelos.

El estilo importa. Los requisitos se enuncian en frases cortas, una condición por
requisito, con identificador propio y sin conjunciones que escondan dos
exigencias en una. Las palabras que expresan obligación deben usarse con un
significado acordado y constante, y conviene evitar los términos que no admiten
prueba: «adecuado», «amigable», «eficiente», «si es posible».

\subsection*{Análisis}

El análisis transforma los requisitos en modelos que permiten razonar sobre ellos
antes de diseñar. En un enfoque orientado a objetos, el resultado típico es el
\textbf{modelo conceptual} o modelo del dominio: las clases del problema, sus
atributos y las asociaciones entre ellas \cite{larman2006}.

Ese modelo no es un diseño. Recoge los conceptos del dominio tal como existen para
el cliente, sin métodos, sin decisiones de implementación y sin clases técnicas.
Confundir ambas cosas produce un modelo del dominio contaminado de controladores y
gestores, que ya no sirve para hablar con quien conoce el negocio.

Junto al modelo conceptual se emplean los \textbf{diagramas de secuencia del
sistema}, que muestran las operaciones que un actor invoca sobre el sistema visto
como una caja negra, y los \textbf{contratos de operación}, que precisan el efecto
de cada una en términos de precondiciones y postcondiciones.

\subsection*{Validación y gestión de los requisitos}

Validar los requisitos es comprobar que describen el sistema que el cliente quiere,
frente a verificar, que es comprobar que el sistema construido cumple los
requisitos. Las técnicas habituales son las revisiones formales con los
implicados, la construcción de prototipos y la redacción anticipada de los casos
de prueba: un requisito para el que no se sabe escribir una prueba está mal
enunciado.

Los requisitos cambian durante todo el proyecto, así que necesitan gestión: una
línea base acordada, un procedimiento para proponer y aprobar cambios, y
trazabilidad para saber qué se ve afectado por cada uno. Sin trazabilidad, el
coste de un cambio se estima a ojo y siempre por debajo.
